\lesson{3}{Jan 10 2022 Mon (16:54:00)}{The Decision}{Unit 4}

\subsubsection{Types of Judicial Opinions}

When a federal court makes a decision on a case, it publishes an official
document explaining the ruling, called an opinion. In courts with a panel of
judges that rule by majority, such as the Supreme Court, it is possible that not
all of the judges will agree on the final decision. For this reason, multiple
opinions may be published for each case.

The majority opinion reflects the official ruling, the decision to which the
majority of justices agree. The Chief Justice of the Supreme Court, if part of
the majority, will write the majority opinion. If the Chief Justices is not part
of the majority, he or she will assign an associate justice to write the formal
opinion. A justice in the minority will write a minority opinion, also called a
dissent or dissenting opinion. In it, the justice describes why he or she does
not agree with the majority and what the ruling should have been. The other
justices may add their signatures to the majority or minority opinion as
appropriate, though often one or more will write a separate concurring opinion.
A concurring opinion agrees with either the majority or minority opinion, but
for different reasons or while disagreeing with a specific part of the opinion.

A per curiam is a ruling explained by the court as a whole, rather than
individual opinions of all the justices involved. The Supreme Court rarely
writes a per curiam decision.

\subsubsection{Interpretations Differ}

Many things can affect a judge's interpretation of the law. One of the most
important is Stare Decisis, meaning prior decisions on the case or related
situations. The decisions of the lower courts on a specific case factor in to
the judge's decisions, as well as their own court's decisions on related cases
from the past. These prior decisions are called precedents.

The way a judge interprets and applies the laws, precedents, and Constitution to
a particular case can vary, just as you can be flexible or literal in the way
you understand language. "Strict constructionist" describes judges that favor a
literal interpretation of the laws, court decisions, and Constitution.
"Activist" or "loose constructionists" describes judges who have a more flexible
view of the law, seeing it as able to apply to new situations not described or
anticipated by the original law.

Americans debate whether one style of interpretation is more appropriate than
the other. The term "activist judge" has a negative connotation for many,
implying a judge who makes decisions beyond the scope of his or her power.
However, others note that judges who have favored more loose interpretations of
the law are precisely the reason Americans have more protection for civil rights
today.

\subsubsection{Symbolic Speech}

Some of the most famous Supreme Court rulings involve the First Amendment from
the Bill of Rights, which protects the freedoms of religion, speech, press,
assembly, and petition-

Remember that rights are not absolute, so laws and court rulings have affected
the way they are interpreted. The meaning of the term "free speech" has been the
subject of heavy debate. View this slideshow to see how two landmark Supreme
Court cases have affected interpretation of the term "free speech."

\newpage

